\documentclass[10pt,a4paper]{amsart}
\usepackage[margin=19mm]{geometry}
\usepackage{amsmath,amssymb,mathtools}
\usepackage[scr=rsfs,cal=euler]{mathalfa}
\usepackage{xfrac}
\usepackage[unicode,hidelinks,bookmarks=false]{hyperref}
\newcommand{\ie}{i.e.\ }
\newcommand{\cf}{cf.\ }
\newcommand{\resp}{respectively\ }
\newcommand{\Subsection}{\subsection}
\providecommand{\nobreakdash}{\nobreak\hbox{-}}
\setlength{\emergencystretch}{2em}
\multlinegap0pt
\usepackage{enumerate}
\newtheorem{prop}{Proposition}
\newtheorem{cor}{Corollary}
\usepackage{cleveref}
\crefname{prop}{Proposition}{Propositions}
\crefname{cor}{Corollary}{Corollaries}
\title[prop:nonempty]{$7$-located locally $5$-large complexes are aspherical (source label prop:nonempty)}
\author{Katherine Goldman$^{\dag}$, Piotr Przytycki$^{\ddag}$}
\date{Source: arXiv:2509.25674v2 (CC BY 4.0)}
\begin{document}
\maketitle

\noindent\emph{Source and licence.} $7$-located locally $5$-large complexes are aspherical, by Katherine Goldman$^{\dag}$, Piotr Przytycki$^{\ddag}$. Version arXiv:2509.25674v2 (CC BY 4.0),
distributed by arXiv under the Creative Commons Attribution 4.0 International licence,
http://creativecommons.org/licenses/by/4.0/, as recorded in the arXiv licence metadata for this exact version
and shown on the abstract page https://arxiv.org/abs/2509.25674v2. Full licence notice:
\texttt{licenses/arxiv-2509.25674-CC-BY-4.0.txt}.\par\smallskip

\noindent\textit{Editorial scope.} Counted unit: the article's Proposition (source0.tex lines 396--399, source label \texttt{prop:nonempty}): for $n>0$ and any simplex $\sigma$ of $S_n(x)$, the subcomplex $K(\sigma)$ is nonempty. The article's complete original proof of it is delivered with it (source0.tex lines 401--410, one \texttt{proof} environment of 10 lines, adjacent to the statement, with no gap and no deferral). No mathematics is written, repaired or re-derived: the statement and the proof are the article's own lines, in the article's own notation, and no retained line is substituted, re-flowed or omitted, so nothing is removed and the package carries no omission record. Retained uncounted as the source-local context the proof needs: the statement of the diameter proposition that the first case invokes (lines 366--373, source label \texttt{prop: diam of K}). Every cross-reference of every retained line resolves inside this document. No retained line contains a citation, so the wrapper carries no bibliography and no bibliography entry; note that the article's other citations have no bibliography source in the e-print, which is why the counted unit was chosen among the cite-free results. No retained line contains a figure environment, an image-inclusion command or drawing code, so no image is delivered and the package declares no asset dependency.

Editorial formatting only, in addition: an AMS \texttt{amsart} preamble that restates the article's own declarations used by the retained lines, namely the environments \texttt{prop} and \texttt{cor} on separate counters so that the article's own \texttt{\textbackslash Cref} cross-references print the correct environment names, together with the standard \texttt{cleveref} package the article itself loads. The retained environments print in this document as Corollary 1 and Proposition 2 in order of appearance, whereas the article prints the same two items with the numbering of its own full document; the retained lines contain no numbered display.

The complete native source of the article as distributed in the arXiv e-print for this exact version is bundled unchanged as \texttt{sources/186030/source0.tex}: it is byte identical to the e-print file \texttt{main\_revised.tex} except that the pipeline's mechanical concatenation prepends one header comment line naming it. The article uses the AMS \texttt{amsart} class; no publisher class or style file is shipped in the e-print and none is redistributed or used.
\begin{cor} \label{prop: diam of K}
Let $v_1,v_2$ be vertices of $K(v)$.
\begin{enumerate}[(i)]
    \item If $v_1$ and $v_2$ are neighbours, then $K(v_1)$ intersects $ K(v_2)$.
    \item If $v_1$ and $v_2$ are not neighbours, then 
    they have a common neighbour $v'$ in~$K(v)$ such that there is an edge $v_1'v_2'$ with $v_1'$ in $K(v_1v')$ and $v_2'$ in $K(v_2v')$.
\end{enumerate}
\end{cor}

\begin{prop}
\label{prop:nonempty}
   Let $n > 0$, and let $\sigma$ be a simplex of $S_n(x)$. Then $K(\sigma)$ is nonempty.
\end{prop}

\begin{proof}
    Suppose first that $\sigma$ is an edge of $S_n(x)$ with vertices $v_1$ and $v_2$. We may obtain a new complex $X'$ by artificially adding to $X$ a vertex $v$ and a triangle $vv_1v_2$.
    This does not affect local $5$-largeness or $7$-location, so the proposition follows from \Cref{prop: diam of K}(i) applied to $v$ in $X'$.

    Now suppose $\dim(\sigma) > 1$. We fix two distinct vertices $v$ and $y$ of $\sigma$. Let $\sigma'$ be the subsimplex of $\sigma$ spanned on all the vertices except for $y$, and let $e=vy$.
    By induction, 
    we have vertices $v_1$ in $K(\sigma')$ and $v_2$ in $K(e)$.
    If neither $v_1$ nor $v_2$ lie in~$K(\sigma)$, then there is $u\neq v$ in $\sigma'$ that is not a neighbour of $v_2$, and $y$ is not a neighbour of $v_1$. By the $5$-largeness of the link of $v$, the vertex $v_1$ is not a neighbour of $v_2$. Let $v',v'_1,v'_2$ be the vertices from Corollary~\ref{prop: diam of K}(ii). Note that if $v'$ is a neighbour of $y$, then by the $5$-largeness of the link of $v$ it lies in $K(\sigma')$. Thus we can assume that $v'$ is not a neighbour of $y$ and so $W_1=(v, v_2yuv_1v')$ is a $5$-wheel. Since $W_2=(v',v_2v_2'v_1'v_1v)$ is also a $5$-wheel, $(W_1,W_2)$ 
    is a $(5,5)$-dwheel, and so all the vertices of $W_1\cup W_2$ have a common neighbour $z$, which lies in $K(v)$. By the $5$-largeness of the link of $v$, the vertex $z$ lies in $K(\sigma')$, and hence also in $K(\sigma)$.
\end{proof}

\end{document}
